
\begin{figure}[t!]
\centering
\begin{tikzpicture}[
    x=0.78cm,
    y=1cm,
    >=Latex,
    line cap=butt
]

    %------------------------------------------------------------
    % Geometry
    %------------------------------------------------------------
    \def\boxheight{2.8}
    \def\boxwidth{16}

    % End of the first and second boxes
    \def\startend{1.6}
    \def\delayend{3.6}

    % Partitions inside the walk region
    \def\walkA{6.4}
    \def\walkB{7.55}
    \def\walkC{8.70}
    \def\walkD{9.85}

    % Half-length of each small horizontal interval separator
    \def\separatorhalfwidth{0.12}

    % Vertical midpoint of the large box
    \pgfmathsetmacro{\ymid}{\boxheight/2}

    % Centers of the first two boxes
    \pgfmathsetmacro{\startcenter}{\startend/2}
    \pgfmathsetmacro{\delaycenter}{(\startend+\delayend)/2}

    % Centers of the three narrow boxes
    \pgfmathsetmacro{\dcenterA}{(\walkA+\walkB)/2}
    \pgfmathsetmacro{\dcenterB}{(\walkB+\walkC)/2}
    \pgfmathsetmacro{\dcenterC}{(\walkC+\walkD)/2}

    % The walk region is everything after the delay box
    \pgfmathsetmacro{\walkcenter}{(\delayend+\boxwidth)/2}

    %------------------------------------------------------------
    % Styles
    %------------------------------------------------------------
    \tikzset{
        outer box/.style={
            draw=black,
            rounded corners=2mm,
            line width=0.65pt
        },
        partition/.style={
            draw=black,
            line width=0.55pt
        },
        colored bar/.style={
            line width=3pt,
            line cap=butt
        },
        interval separator/.style={
            draw=black,
            line width=0.55pt,
            line cap=round
        },
        small arrow/.style={
            -{Latex[length=1.8mm,width=1.35mm]},
            line width=0.9pt,
            black
        },
        upward arrow/.style={
            -{Latex[length=2.7mm,width=2.1mm]},
            line width=1.5pt,
            green!65!black,
            line cap=butt
        },
        downward arrow/.style={
            -{Latex[length=2.7mm,width=2.1mm]},
            line width=1.5pt,
            red!70,
            line cap=butt
        },
        lower label/.style={
            anchor=north,
            font=\large\itshape
        },
        set label/.style={
            font=\Large
        }
    }

    %------------------------------------------------------------
    % Clip all interior elements to the rounded outer box
    %------------------------------------------------------------
    \begin{scope}
        \clip[rounded corners=2mm]
            (0,0) rectangle (\boxwidth,\boxheight);

        %--------------------------------------------------------
        % Exact top-to-bottom partition lines
        %--------------------------------------------------------
        \foreach \x in {
            \startend,
            \delayend,
            \walkA,
            \walkB,
            \walkC,
            \walkD
        }{
            \draw[partition]
                (\x,0)
                --
                (\x,\boxheight);
        }

        %--------------------------------------------------------
        % First box: unequal interval lengths
        %--------------------------------------------------------
        \draw[colored bar,green!65!black]
            (\startcenter,0)
            --
            (\startcenter,0.55);

        \draw[colored bar,red!75]
            (\startcenter,0.55)
            --
            (\startcenter,1.90);

        \draw[colored bar,blue!60]
            (\startcenter,1.90)
            --
            (\startcenter,\boxheight);

        % Horizontal separators between the three intervals
        \foreach \y in {0.55,1.90}{
            \draw[interval separator]
                ({\startcenter-\separatorhalfwidth},\y)
                --
                ({\startcenter+\separatorhalfwidth},\y);
        }

        %--------------------------------------------------------
        % Second box: six equal intervals
        %--------------------------------------------------------
        \foreach \col [count=\i from 0] in {
            red!75,
            orange!85!black,
            cyan!65!black,
            green!60!black,
            magenta!65,
            blue!60
        }{
            \pgfmathsetmacro{\segmentbottom}{\i*\boxheight/6}
            \pgfmathsetmacro{\segmenttop}{(\i+1)*\boxheight/6}

            \draw[colored bar,\col]
                (\delaycenter,\segmentbottom)
                --
                (\delaycenter,\segmenttop);
        }

        % Horizontal separators between the six equal intervals
        \foreach \i in {1,...,5}{
            \pgfmathsetmacro{\separatorheight}
                {\i*\boxheight/6}

            \draw[interval separator]
                ({\delaycenter-\separatorhalfwidth},
                 \separatorheight)
                --
                ({\delaycenter+\separatorhalfwidth},
                 \separatorheight);
        }

        %--------------------------------------------------------
        % Three narrow boxes:
        % each is divided into equal upper and lower intervals
        %--------------------------------------------------------
        \foreach \x in {
            \dcenterA,
            \dcenterB,
            \dcenterC
        }{
            % Upper interval
            \draw[upward arrow]
                (\x,\ymid)
                --
                (\x,\boxheight);

            % Lower interval
            \draw[downward arrow]
                (\x,\ymid)
                --
                (\x,0);

            % Horizontal separator between the two intervals
            \draw[interval separator]
                ({\x-\separatorhalfwidth},\ymid)
                --
                ({\x+\separatorhalfwidth},\ymid);
        }

        %--------------------------------------------------------
        % Small left-to-right arrows
        %--------------------------------------------------------

        % First box
        \draw[small arrow]
            (0.32,1.23)
            --
            (0.59,1.23);

        % Second box
        \draw[small arrow]
            (2.08,2.15)
            --
            (2.35,2.15);

        % First narrow box
        \draw[small arrow]
            (6.57,2.02)
            --
            (6.82,2.02);

        % Second narrow box
        \draw[small arrow]
            (7.70,0.70)
            --
            (7.95,0.70);

        % Third narrow box
        \draw[small arrow]
            (8.85,0.95)
            --
            (9.10,0.95);
    \end{scope}

    %------------------------------------------------------------
    % Draw the outer boundary last
    %------------------------------------------------------------
    \draw[outer box]
        (0,0)
        rectangle
        (\boxwidth,\boxheight);

    %------------------------------------------------------------
    % External labels
    %------------------------------------------------------------
\node[set label,overlay]
    at (-0.72,\ymid)
    {$[C]$};

    \node[set label]
        at ({\boxwidth/2},{\boxheight+0.62})
        {$[P]$};

    % All lower labels share the same baseline
    \node[lower label]
        at (\startcenter,-0.1)
        {\pstart};

    \node[lower label]
        at (\delaycenter,-0.1)
        {\pdelay};

    \node[lower label]
        at (\dcenterA,-0.1)
        {$d+2$};

    \node[lower label]
        at (\walkcenter,-0.1)
        {\pwalk};

\end{tikzpicture}

\caption{An illustration of the random walk $W(x)$ in~\Cref{def:W(x)} when $m=3$ and $D=6$. The figure shows how the indices of $[P]$ are partitioned into $(\pstart,\pdelay,\pwalk)$. The intervals inside each box shows how $[C]$ is partitioned into different intervals. The tiny arrows in each box shows the value of a single point $x \in [C]^P$ on each relevant index. The walk $W(x)$ thus first considers $x_{\pstart}$ and so chooses starting state $2$. The delay $d$ is set to $5$ given $x_{\pdelay}$. Finally, the walk waits until time $d$ and then decide its next move based on the index $d+2$ and so on; in this particular example, the walk goes to $3$ next, then $2$, and then $1$. }
\label{fig:W(x)}
\end{figure}